% Rantai kepercayaan sertifikat X.509.
% Dirujuk dari section-03.tex dengan % Rantai kepercayaan sertifikat X.509.
% Dirujuk dari section-03.tex dengan % Rantai kepercayaan sertifikat X.509.
% Dirujuk dari section-03.tex dengan % Rantai kepercayaan sertifikat X.509.
% Dirujuk dari section-03.tex dengan \input{figures/rantai-sertifikat}.
\begin{figure}[htbp]
    \centering
    \begin{tikzpicture}[
        kotak/.style={draw, rounded corners, align=center, font=\small, inner sep=5pt,
                      minimum width=6.4cm, minimum height=1.0cm},
        panah/.style={-{Latex[length=2.4mm]}, thick},
        node distance=1.0cm,
    ]
        \node[kotak, fill=gray!12] (root) {\textbf{Root CA}\\ \scriptsize kunci publik tertanam di peramban dan sistem operasi};
        \node[kotak, fill=blue!6, below=of root] (inter) {\textbf{Intermediate CA}\\ \scriptsize sertifikatnya ditandatangani Root CA};
        \node[kotak, fill=green!10, below=of inter] (server) {\textbf{Sertifikat Server}\\ \scriptsize ditandatangani Intermediate CA};

        \draw[panah] (root) -- node[right, font=\scriptsize, align=left] {menandatangani} (inter);
        \draw[panah] (inter) -- node[right, font=\scriptsize, align=left] {menandatangani} (server);
    \end{tikzpicture}
    \caption{Rantai kepercayaan sertifikat X.509: kepercayaan pada Root CA diteruskan ke sertifikat server melalui Intermediate CA}
    \label{fig:rantai-sertifikat}
\end{figure}
.
\begin{figure}[htbp]
    \centering
    \begin{tikzpicture}[
        kotak/.style={draw, rounded corners, align=center, font=\small, inner sep=5pt,
                      minimum width=6.4cm, minimum height=1.0cm},
        panah/.style={-{Latex[length=2.4mm]}, thick},
        node distance=1.0cm,
    ]
        \node[kotak, fill=gray!12] (root) {\textbf{Root CA}\\ \scriptsize kunci publik tertanam di peramban dan sistem operasi};
        \node[kotak, fill=blue!6, below=of root] (inter) {\textbf{Intermediate CA}\\ \scriptsize sertifikatnya ditandatangani Root CA};
        \node[kotak, fill=green!10, below=of inter] (server) {\textbf{Sertifikat Server}\\ \scriptsize ditandatangani Intermediate CA};

        \draw[panah] (root) -- node[right, font=\scriptsize, align=left] {menandatangani} (inter);
        \draw[panah] (inter) -- node[right, font=\scriptsize, align=left] {menandatangani} (server);
    \end{tikzpicture}
    \caption{Rantai kepercayaan sertifikat X.509: kepercayaan pada Root CA diteruskan ke sertifikat server melalui Intermediate CA}
    \label{fig:rantai-sertifikat}
\end{figure}
.
\begin{figure}[htbp]
    \centering
    \begin{tikzpicture}[
        kotak/.style={draw, rounded corners, align=center, font=\small, inner sep=5pt,
                      minimum width=6.4cm, minimum height=1.0cm},
        panah/.style={-{Latex[length=2.4mm]}, thick},
        node distance=1.0cm,
    ]
        \node[kotak, fill=gray!12] (root) {\textbf{Root CA}\\ \scriptsize kunci publik tertanam di peramban dan sistem operasi};
        \node[kotak, fill=blue!6, below=of root] (inter) {\textbf{Intermediate CA}\\ \scriptsize sertifikatnya ditandatangani Root CA};
        \node[kotak, fill=green!10, below=of inter] (server) {\textbf{Sertifikat Server}\\ \scriptsize ditandatangani Intermediate CA};

        \draw[panah] (root) -- node[right, font=\scriptsize, align=left] {menandatangani} (inter);
        \draw[panah] (inter) -- node[right, font=\scriptsize, align=left] {menandatangani} (server);
    \end{tikzpicture}
    \caption{Rantai kepercayaan sertifikat X.509: kepercayaan pada Root CA diteruskan ke sertifikat server melalui Intermediate CA}
    \label{fig:rantai-sertifikat}
\end{figure}
.
\begin{figure}[htbp]
    \centering
    \begin{tikzpicture}[
        kotak/.style={draw, rounded corners, align=center, font=\small, inner sep=5pt,
                      minimum width=6.4cm, minimum height=1.0cm},
        panah/.style={-{Latex[length=2.4mm]}, thick},
        node distance=1.0cm,
    ]
        \node[kotak, fill=gray!12] (root) {\textbf{Root CA}\\ \scriptsize kunci publik tertanam di peramban dan sistem operasi};
        \node[kotak, fill=blue!6, below=of root] (inter) {\textbf{Intermediate CA}\\ \scriptsize sertifikatnya ditandatangani Root CA};
        \node[kotak, fill=green!10, below=of inter] (server) {\textbf{Sertifikat Server}\\ \scriptsize ditandatangani Intermediate CA};

        \draw[panah] (root) -- node[right, font=\scriptsize, align=left] {menandatangani} (inter);
        \draw[panah] (inter) -- node[right, font=\scriptsize, align=left] {menandatangani} (server);
    \end{tikzpicture}
    \caption{Rantai kepercayaan sertifikat X.509: kepercayaan pada Root CA diteruskan ke sertifikat server melalui Intermediate CA}
    \label{fig:rantai-sertifikat}
\end{figure}
